% FILE: 01_research_question.tex
% PROJECT: sources-papers
% ROLE: academic-paper-classification-and-evidence-grounding

\section{Research Question}
\label{sec:sources-papers-rq}

\subsection{Problem Statement}
\label{sec:sources-papers-rq-problem}

The process mining literature has, over three decades, accumulated a corpus of
formal results --- soundness criteria for Workflow nets, discovery algorithms
for Petri nets and process trees, conformance metrics bounded in $[0,1]$,
object-centric log specifications, declarative constraint templates --- each
stated in academic notation. None of these results, by themselves, prevent an
implementation from violating the invariants they define. A codebase may claim
to produce a ``sound'' WF-net while its return type is indistinguishable from
an unsound one. A system may claim ``alignment-based fitness $\geq 0.95$''
while never executing an alignment algorithm at all.

The problem this project addresses is the \emph{formalization gap}: the
distance between a published invariant and a compiler-enforced guarantee. The
paper corpus in \texttt{sources/papers} exists precisely to close this gap by
deriving, from each paper's formal objects and algorithms, a precise set of
obligations that must exist in the Rust type system and the wasm4pm execution
engine before any downstream claim can be treated as evidence rather than
assertion.

A secondary problem is the \emph{adversarial gap}: canonical process mining
algorithms structurally fail when event logs are manipulated by Sybil
attackers. Case ID forgery causes state-space explosion and yields incorrectly
permissive discovered models. Non-deterministic reality receipts cause
canonical discovery algorithms to misinterpret adversarial traces as noise or
novel process deviants. No existing published algorithm provides
Byzantine-fault-tolerant discovery; this gap is formally documented in
\texttt{adversarial-canon.md}.

\subsection{Old-Regime Assumption Challenged}
\label{sec:sources-papers-rq-old-regime}

The old-regime assumption challenged by this project is that \emph{academic
citation is sufficient for claim substantiation}. Under the old regime, a
system vendor may cite van der Aalst (1998) to support a claim of WF-net
soundness without demonstrating that the implementation enforces the
three-condition soundness criterion (option to complete, proper completion,
no dead transitions) at any verifiable boundary. The citation acts as social
proof, not as a compiler-enforced invariant.

A related old-regime assumption is that \emph{PM4Py coverage implies wasm4pm
coverage}. Because PM4Py implements Alpha Miner, Inductive Miner,
alignment-based conformance, and dozens of other algorithms in Python, a
system that wraps or mirrors PM4Py's API surface is presumed to inherit those
capabilities. This project falsifies that assumption directly: all 29
algorithms enumerated in \texttt{PAPER\_TO\_EXECUTION\_LAW.md} carry
\texttt{NO} in the wasm4pm coverage column as of 2026-05-31. PM4Py coverage
is documented as an \emph{oracle reference}, not as an implementation
substitute.

A third old-regime assumption is that \emph{deterministic event logs are the
normal case}. The adversarial-canon analysis demonstrates that this assumption
is a security vulnerability, not a reasonable operating condition.

\subsection{Hypothesis}
\label{sec:sources-papers-rq-hypothesis}

The central hypothesis of this project is:

\begin{quotation}
\noindent\textbf{Hypothesis SP-1.} Every formal invariant stated in a canonical
process mining paper can be translated into a zero-cost Rust type law such that
programs violating the invariant fail to compile, and such that the compilation
failure message names the invariant and cites the originating paper. This
translation is possible for all \emph{structural} invariants; invariants
requiring dynamic graph reachability computation are translatable to a named
graduation boundary that defers computation to the execution engine without
abandoning compile-time safety.
\end{quotation}

A corollary hypothesis governs the adversarial domain:

\begin{quotation}
\noindent\textbf{Hypothesis SP-2.} Byzantine-fault-tolerant process discovery
requires replacing case-ID correlation with BLAKE3 cryptographic receipt
verification. Any discovery algorithm that accepts an event log without first
verifying receipt provenance is structurally vulnerable to Sybil attack,
regardless of algorithmic sophistication.
\end{quotation}

\subsection{Success Criteria}
\label{sec:sources-papers-rq-success}

Success for this project is defined by the following measurable criteria,
derived from the ALIVE gate assessment (2026-05-31):

\begin{enumerate}
  \item \textbf{File count}: \texttt{sources/papers} $\geq 7$ files. Actual:
    15. \emph{MET.}
  \item \textbf{Classification completeness}: Every paper in the 81-paper
    corpus carries one of the five coverage statuses with a stated reason.
    \emph{MET} (documented in \texttt{PAPER\_TO\_TYPE\_LAW.md}).
  \item \textbf{Type-law coverage}: 18 papers fully covered with named Rust
    types in \texttt{wasm4pm-compat}. \emph{MET.}
  \item \textbf{Obligation backlog}: All missing types named explicitly, not
    left implicit. 22 named missing types and 14 named missing output shapes
    on record. \emph{MET.}
  \item \textbf{Board-claim traceability}: Each Tier-1 board claim carries an
    evidence path naming specific source files and compile-fail fixtures.
    \emph{MET} for Tier-1 structural completeness claims.
  \item \textbf{Receipt}: PAPER\_CANON\_RECEIPT issued with zero gaps
    remaining. \emph{MET.}
  \item \textbf{Adversarial coverage}: Byzantine threat model documented with
    named attack vectors (Sybil, non-deterministic reality receipts) and
    corresponding required defences (BLAKE3, typestate progression).
    \emph{MET} at documentation level; implementation evidence absent (PARTIAL).
\end{enumerate}
